% ----------------------------------------------------------
% Conclusão
% ----------------------------------------------------------
\chapter[Conclusão]{Conclusão}
%\addcontentsline{toc}{chapter}{Conclusão}
% ---

Este trabalho aplicou análise de redes complexas a dados históricos olímpicos, revelando padrões estruturais de competição e hierarquias de performance através de métricas de centralidade e detecção de comunidades. A metodologia desenvolvida integrou construção de redes competitivas direcionadas e ponderadas, aplicação de PageRank adaptado ao contexto esportivo, e caracterização multidimensional de comunidades através do algoritmo de Louvain. Os resultados obtidos confirmam a hipótese central de que análise de redes revela estruturas competitivas latentes não capturadas por métricas tradicionais de contagem de medalhas.

\section{Síntese dos Achados Principais}

A análise de 9.350 atletas medalhistas distribuídos em seis modalidades olímpicas ao longo de 125 anos (1896-2021, incluindo Tokyo 2020) revelou diferenças estruturais sistemáticas entre tipos de competição esportiva. Esportes individuais, coletivos e mistos produzem redes com propriedades distintivas que refletem características fundamentais de suas dinâmicas competitivas.

\textbf{Diferenças Estruturais entre Modalidades:} Swimming apresentou maior modularidade (0.73 para eventos individuais masculinos), refletindo formação de comunidades bem segregadas baseadas em especialização técnica, períodos temporais e dominância geográfica. Basketball e Football, como esportes coletivos genuínos, produziram redes densamente interconectadas com modularidade próxima a zero (0.003 e 0.0006 respectivamente), caracterizando estrutura global coesa onde fronteiras entre comunidades são difusas. Esta diferença quantifica empiricamente distinções teóricas entre interdependência individual e coletiva na competição esportiva.

\textbf{Comunidades Refletem Múltiplos Fatores:} As 166 comunidades detectadas não correspondem a agrupamentos simples por período temporal ou nacionalidade, mas emergem de combinações complexas de fatores. A caracterização multidimensional revelou que comunidades em natação apresentam diversidade temporal moderada (entropia média 1.60), enquanto esportes coletivos exibem diversidade temporal superior (Basketball: 2.46, Football: 2.53), refletindo que comunidades em esportes coletivos agregam atletas de múltiplas décadas devido à natureza densamente conectada destas redes. Modalidades de combate individual (Judo, Boxing) demonstraram segregação estrutural ainda mais acentuada (modularidade média 0.81), com fronteiras delimitadas primariamente por categorias de peso. Esta heterogeneidade indica que algoritmos de detecção de comunidades capturam estruturas emergentes não redutíveis a características isoladas.

\textbf{Métricas de Rede Revelam Padrões Ocultos:} PageRank identificou hierarquias de importância estrutural distintas de rankings por contagem de medalhas. O ranking global é dominado por boxeadoras (16 das 20 primeiras posições), com Claressa Shields (USA) liderando com PageRank de 0.0822, valor 3.2× superior ao de Michael Phelps (0.0255, 11º global), apesar deste possuir 28 medalhas olímpicas contra 2 de Shields. Esta inversão evidencia limitação metodológica crítica: PageRank quantifica importância relativa dentro da topologia de cada rede, não performance absoluta entre redes de tamanhos distintos. Redes pequenas (Boxing F: 41 atletas) concentram PageRank em poucos nós, enquanto redes massivas (Football M: 1.275 atletas) diluem mesmo campeões históricos. Esta distinção exemplifica como métricas de rede capturam propriedades estruturais invisíveis em análises convencionais, mas exigem interpretação contextual criteriosa.

\textbf{Atletas-Ponte Conectam Comunidades:} A identificação de atletas com alta betweenness centrality revelou indivíduos estruturalmente críticos que conectam diferentes regiões das redes. Em natação, 10.2\% dos atletas atuam como pontes, proporção superior aos esportes coletivos (5.3-7.1\%). Estes atletas frequentemente apresentam longevidade olímpica ou versatilidade técnica, posicionando-se como conectores entre eras competitivas ou especializações distintas. Sua importância estrutural transcende performance individual, evidenciando papel sistêmico não capturado por métricas tradicionais.

\section{Contribuições Metodológicas}

Este trabalho entrega três contribuições principais para a literatura de análise de redes aplicada ao esporte:

\textbf{Caracterização Multidimensional de Comunidades:} Desenvolvemos framework de análise que vai além de mera detecção de fronteiras comunitárias, integrando métricas temporais (entropia de Shannon), geográficas (diversidade de NOCs) e de performance (coeficiente de variação de PageRank). Esta abordagem revelou que comunidades em esportes individuais exibem alta segregação temporal e geográfica, enquanto esportes coletivos apresentam maior homogeneidade estrutural.

\textbf{Análise Comparativa Multi-Esporte:} A aplicação da metodologia a seis modalidades olímpicas com diferentes tipos de interdependência competitiva (individual puro, coletivo, misto, combate individual) abrangendo 125 anos de história olímpica permite identificação de padrões estruturais sistemáticos relacionados ao formato competitivo. A diferença de modularidade entre Swimming (0.73), esportes de combate (Judo/Boxing: ~0.81), e esportes coletivos genuínos (Basketball/Football: ~0.003) quantifica empiricamente distinções teóricas entre tipos de competição, validando aplicabilidade de análise de redes para caracterização de estruturas competitivas diversas.

\textbf{Dashboard Interativo para Exploração:} A implementação de interface web interativa oferece ferramenta prática para exploração dos resultados através de múltiplas perspectivas analíticas. O dashboard integra análise de centralidades, visualização de comunidades, análise temporal e comparações entre esportes, facilitando identificação de padrões e validação de hipóteses sobre estruturas competitivas. Esta contribuição demonstra viabilidade de ferramentas acessíveis para exploração de redes complexas por audiências não-técnicas.

\section{Limitações do Estudo}

Este trabalho apresenta limitações que devem ser consideradas na interpretação dos resultados:

\textbf{Limitações de Dados:}

\textit{Viés de survivorship histórico:} Registros de atletas pré-1960 são fragmentados, potencialmente sub-representando medalhistas de eras anteriores. Dados incompletos sobre altura, peso e idade em períodos iniciais dos jogos limitam análises que correlacionam atributos físicos com posições estruturais.

\textit{Cobertura temporal:} Análise limitada a Tokyo 2020, não incluindo Paris 2024. A inclusão de edições futuras permitiria análise de tendências mais recentes em profissionalização e globalização.

\textit{Escopo de modalidades:} Foco em seis esportes (~7,1\% das modalidades olímpicas), impossibilitando generalização para todas as 84 modalidades. Esportes com características distintas (eventos cronometrados puros, modalidades artísticas) podem produzir propriedades estruturais não capturadas.

\textbf{Limitações Metodológicas:}

\textit{Sistema de ponderação exploratório:} Escolha do sistema 5-3-2 carece de fundamentação teórica sólida, sendo validada apenas empiricamente via análise de robustez. Embora validação tenha demonstrado concordância estrutural entre sistemas alternativos, justificativa teórica para pesos específicos permanece aberta.

\textit{Diluição estrutural de PageRank:} Esportes coletivos com múltiplos medalhistas por evento diluem PageRank artificialmente, dificultando comparações diretas com esportes individuais. A formulação recursiva da métrica distribui importância proporcionalmente ao tamanho da rede, favorecendo estruturalmente competições com menor número de participantes.

\textit{Perda de informação direcional:} Algoritmo de Louvain requer simetrização de grafos direcionados, descartando hierarquias de dominância durante detecção de comunidades. Embora validação com Infomap (que preserva direção) tenha demonstrado concordância estrutural de 85\% (NMI), informação direcional potencialmente relevante foi perdida na análise comunitária.

\textit{Validação estatística de comunidades:} Validação formal através de null models (comparação com redes aleatórias preservando distribuição de grau) não foi implementada. Magnitude das diferenças de modularidade (Swimming 0.73 vs Basketball/Football ~0.003, razão de 243×) suporta conclusões qualitativas, mas significância estatística formal não foi estabelecida.

\textbf{Limitações de Escopo:}

\textit{Ausência de dados contextuais:} Análise não incorpora investimento esportivo, políticas de inclusão, PIB dos países ou mudanças regulatórias que impactam estrutura competitiva. Fatores socioeconômicos e políticos que explicam posições estruturais permanecem inexplorados.

\textit{Foco em medalhistas:} Exclusão de atletas não-medalhistas ignora estrutura competitiva completa, capturando apenas topo da hierarquia. Padrões de participação, distribuição geográfica e diversidade em níveis competitivos inferiores não foram analisados.

\section{Trabalhos Futuros}

Os resultados obtidos e limitações identificadas sugerem múltiplas direções promissoras para extensão desta pesquisa.

\textbf{Expansão para Mais Modalidades:} A aplicação da metodologia desenvolvida a conjunto mais abrangente de esportes olímpicos permitiria validação sistemática dos achados e identificação de propriedades estruturais gerais versus específicas de modalidade. Particular interesse reside em esportes de combate com categorias adicionais (taekwondo, luta olímpica), modalidades cronometradas sem confronto direto (ciclismo, remo), eventos artísticos com avaliação subjetiva (ginástica artística, nado sincronizado) e esportes de raquete (tênis, badminton), que apresentam características estruturais potencialmente distintas das seis modalidades analisadas (Swimming, Basketball, Football, Athletics, Judo, Boxing).

\textbf{Análise de Evolução Topológica de Redes:} Embora o presente trabalho tenha caracterizado dinâmicas temporais de dominância e diversidade através de entropia de Shannon por era olímpica, a construção de sequências temporais de redes (snapshots por década) revelaria evolução da topologia das redes ao longo da história. Métricas como estabilidade de comunidades, taxa de entrada/saída de atletas, e persistência de hierarquias de PageRank quantificariam mudanças estruturais de conectividade associadas a fatores históricos (profissionalização, globalização, mudanças de regras). Modelos de redes temporais permitiriam previsão de tendências futuras em estruturas competitivas.

\textbf{Comparação com Rankings Oficiais:} A validação sistemática de rankings produzidos por PageRank contra sistemas de classificação oficiais (rankings mundiais, pontuações olímpicas, índices de performance) quantificaria convergência ou divergência entre métricas estruturais e avaliações estabelecidas. Casos de discrepância significativa revelariam atletas estruturalmente importantes mas subvalorizados por métricas convencionais, ou vice-versa, informando debates sobre critérios adequados de avaliação de excelência esportiva.

\textbf{Análise de Gênero Sistematizada:} Embora o presente trabalho tenha segregado redes por gênero, análise comparativa sistemática de propriedades estruturais entre competições masculinas e femininas permanece inexplorada. Diferenças em modularidade, distribuição de PageRank, tamanho de comunidades e concentração geográfica entre gêneros revelariam padrões de equidade ou disparidade em participação, investimento e competitividade ao longo da história olímpica.

\textbf{Análise de Redes Multicamada com Contexto Socioeconômico:} Construir redes multicamada integrando camada competitiva (atletas conectados por confrontos) com camada socioeconômica (países conectados por similaridade de PIB per capita e investimento esportivo). Hipótese: PageRank médio de comunidades nacionais correlaciona positivamente com investimento per capita em esporte (r > 0.6), mas saturação ocorre acima de threshold (ex: USD 50/habitante/ano), testável via regressão não-linear. Dados da UNESCO e Banco Mundial permitiriam quantificar retorno marginal decrescente de investimento esportivo sobre dominância competitiva.

O presente trabalho estabeleceu fundamentos metodológicos sólidos para aplicação de análise de redes complexas ao contexto olímpico, demonstrando viabilidade e valor analítico da abordagem. As direções futuras propostas expandiriam compreensão de estruturas competitivas esportivas, contribuindo tanto para avanço teórico em ciência de redes quanto para aplicações práticas em gestão esportiva e políticas de desenvolvimento atlético.